\subsection{Fourier Approximation}

\begin{perspective} \label{per:fourier-plancherel-roadmap-tao} \hlt{Fourier and Plancherel Roadmap}\\
Path: orthonormal characters, Fejer density, best $L^2$ projection, then Parseval/Plancherel.\\
\hlt{Key Insight}: Fourier theory is finite-dimensional orthogonal projection plus an approximation theorem.
\end{perspective}

\begin{definition} \label{def:periodic-convolution-and-fejer-kernel} \hlt{Periodic Convolution and Fejer Kernels}\\
For $f,g\in C(\R/\Z;\C)$, define their \hlt{periodic convolution} as $(f*g)(x)\defeq\int_0^1 f(y)g(x-y)\,dy$.\\
For $N\geq1$, the \hlt{Fejer kernel} is $F_N(x)\defeq \sum\limits_{n=-N}^{N}\left(1-\frac{\abs{n}}{N}\right)e_n(x) =\frac{1}{N}\abs{\sum\limits_{n=0}^{N-1}e_n(x)}^2$.
\end{definition}

\begin{figure}[H]
  \centering
  \begin{tikzpicture}[
  >=Stealth,
  every node/.style={font=\scriptsize},
  x=5.8cm,
  y=0.48cm,
  axis/.style={->, black!55, line width=0.45pt},
  guide/.style={black!35, dashed, line width=0.45pt},
  broad/.style={black!55, dashed, line width=0.75pt},
  sharp/.style={black!85, line width=0.9pt},
  window/.style={fill=gray!12, draw=none}
]
  \draw[window] (-0.10,0) rectangle (0.10,5.65);

  \draw[axis] (-0.52,0) -- (0.55,0) node[right] {$x$};
  \draw[axis] (0,-0.18) -- (0,5.86);

  \draw[guide] (-0.5,0) -- (-0.5,5.65);
  \draw[guide] (0.5,0) -- (0.5,5.65);
  \node[below] at (-0.5,0) {$-\frac12$};
  \node[below] at (0,0) {$0$};
  \node[below] at (0.5,0) {$\frac12$};

  \draw[broad, domain=-0.5:0.5, samples=180, smooth]
    plot (\x,{2.15*exp(-42*\x*\x)});
  \draw[sharp, domain=-0.5:0.5, samples=220, smooth]
    plot (\x,{5.15*exp(-360*\x*\x)});

  \draw[guide] (-0.10,0) -- (-0.10,5.4);
  \draw[guide] (0.10,0) -- (0.10,5.4);
  \draw[<->, black!55] (-0.10,-0.46) -- node[below] {$2\delta$} (0.10,-0.46);

  \node[black!55] at (0.26,1.38) {$F_M$};
  \node[black!85] at (0.12,4.55) {$F_N$};
\end{tikzpicture}

  \caption{The Fejer kernels are nonnegative trigonometric polynomial approximate identities on $\R/\Z$.}
  \label{fig:fejer-kernel-periodic-approx-identity}
\end{figure}

\begin{theorem} \label{thm:fejer-periodic-approximation-identity} \hlt{Fejer Approximation on $\R/\Z$}\\
Let $F_N$ be the period-one Fejer kernel, $N\geq1$. Then:
\begin{enumerate}[label=(\roman*)]
\setlength{\itemsep}{0pt}
\item $F_N$ is a nonnegative trigonometric polynomial and $\int_0^1 F_N(x)\,dx=1$.
\item For each $0<\delta<1/2$, $\sup_{\delta\leq x\leq1-\delta}F_N(x)\to0$.
\item If $f\in C(\R/\Z;\C)$, then $f*F_N$ is a trigonometric polynomial for each $N$.
\item If $f\in C(\R/\Z;\C)$, then $f*F_N\rightrightarrows f$ on $\R/\Z$.
\end{enumerate}
\begin{proof}
\begin{enumerate}[label=(\roman*)]
\setlength{\itemsep}{0pt}
\item By definition, $F_N=\sum\limits_{n=-N}^{N}\left(1-\frac{\abs{n}}{N}\right)e_n$ is trigonometric, and $F_N=\frac{1}{N}\abs{\sum\limits_{n=0}^{N-1}e_n}^2\geq0$.\\
Since $\int_0^1e_n(x)\,dx=0$ for $n\neq0$, the $e_0$ coefficient gives $\int_0^1F_N(x)\,dx=1$.
\item For $x\notin\Z$, $F_N(x)=\frac{\sin^2(\pi Nx)}{N\sin^2(\pi x)}$. If $m_\delta\defeq\min_{\delta\leq x\leq1-\delta}\sin^2(\pi x)>0$, then $0\leq F_N(x)\leq\frac{1}{Nm_\delta}$ on $[\delta,1-\delta]$.
\item Since $f*e_n=\widehat f_{\mathbb{T}}(n)e_n$ and $F_N$ is a finite linear combination of the $e_n$, $f*F_N$ is a trigonometric polynomial.
\item Write $(f*F_N)(x)=\int_0^1f(x-t)F_N(t)\,dt$. Let $M=\norm{f}_\infty$ and $\varepsilon>0$.\\
Uniform continuity gives $0<\delta<1/2$ such that $\abs{f(x-t)-f(x)}<\varepsilon$ for all $x$ and $t\in[0,\delta]\cup[1-\delta,1]$.\\
By (i), $F_N\geq0$ and $\int_0^1F_N(t)\,dt=1$; hence $\abs{(f*F_N)(x)-f(x)}\leq\varepsilon+2M\sup_{\delta\leq t\leq1-\delta}F_N(t)$.\\
By (ii), $\limsup\limits_{N\to\infty}\norm{f*F_N-f}_\infty\leq\varepsilon$. Since $\varepsilon$ is arbitrary, $f*F_N\rightrightarrows f$.
\end{enumerate}
\hlt{Structural Commentary}: Fejer kernels are positive approximate identities.\\
\hlt{Key Insight}: Normalized mass concentrating at $0$ forces uniform convergence.
\end{proof}
\end{theorem}

\begin{perspective} \label{per:fourier-coefficients-as-projections} \hlt{Fourier Coefficients as Projections}\\
Fourier coefficients are projections onto orthonormal directions, modeled finitely by Theorems \ref{thm:orthogonal-direct-sum-decomposition}, \ref{thm:orthonormal-basis-coordinates}, and \ref{thm:bessel-inequality-finite-orthonormal-list}; infinite directions make convergence separate.
\end{perspective}

\begin{theorem} \label{thm:orthonormal-fourier-best-approximation-bessel} \hlt{Fourier Coefficients Give Best Mean-Square Approximation}\\
Let $\langle u,v\rangle\defeq\int_a^b u\overline v$ and $\norm{u}_2^2\defeq\langle u,u\rangle$. Let $(\phi_m)$ be orthonormal and let $f$ satisfy $\norm{f}_2<\infty$.\\
For $n\geq1$, set $c_m\defeq\langle f,\phi_m\rangle$, $s_n\defeq\sum\limits_{m=1}^{n}c_m\phi_m$, and $t_n\defeq\sum\limits_{m=1}^{n}\gamma_m\phi_m$. Then:
\begin{enumerate}[label=(\roman*)]
\setlength{\itemsep}{0pt}
\item $\norm{f-s_n}_2\leq\norm{f-t_n}_2$, with equality iff $\gamma_m=c_m$ for $m=1,\ldots,n$.
\item $\sum\limits_{m=1}^{\infty}\abs{c_m}^2\leq\norm{f}_2^2$, hence $c_m\to0$.
\end{enumerate}
\begin{proof}
\begin{enumerate}[label=(\roman*)]
\setlength{\itemsep}{0pt}
\item By orthonormality, $\norm{f-t_n}_2^2=\norm{f}_2^2-2\operatorname{Re}\sum\limits_{m=1}^{n}c_m\overline{\gamma_m}+\sum\limits_{m=1}^{n}\abs{\gamma_m}^2=\norm{f}_2^2-\sum\limits_{m=1}^{n}\abs{c_m}^2+\sum\limits_{m=1}^{n}\abs{\gamma_m-c_m}^2$. Thus $\norm{f-t_n}_2^2\geq\norm{f-s_n}_2^2$, with equality iff $\gamma_m=c_m$ for all $m\leq n$.
\item Taking $\gamma_m=c_m$ in the identity from (i) gives $\norm{f-s_n}_2^2=\norm{f}_2^2-\sum\limits_{m=1}^{n}\abs{c_m}^2\geq0$.\\
Hence $\sum\limits_{m=1}^{n}\abs{c_m}^2=\norm{s_n}_2^2\leq\norm{f}_2^2$ for all $n$. Letting $n\to\infty$ gives Bessel's inequality.\\
Since $(\abs{c_m}^2)$ is summable, $\abs{c_m}^2\to0$, hence $c_m\to0$.
\end{enumerate}
\hlt{Structural Commentary}: Orthogonality diagonalizes the finite error.\\
\hlt{Key Insight}: Fourier coefficients are projection coordinates.
\end{proof}
\end{theorem}

\begin{theorem} \label{thm:trigonometric-fourier-dirichlet-localization} \hlt{Dirichlet Kernel and Local Fourier Convergence}\\
Let $f$ be Riemann integrable and $2\pi$-periodic, with $\widehat{f}(m)\defeq\frac{1}{2\pi}\int_{-\pi}^{\pi}f(t)e^{-imt}\,dt$.\\
For $N\geq0$, set $s_N(f;x)\defeq\sum\limits_{m=-N}^{N}\widehat{f}(m)e^{imx}$ and $D_N(t)\defeq\sum\limits_{m=-N}^{N}e^{imt}$. Then:
\begin{enumerate}[label=(\roman*)]
\setlength{\itemsep}{0pt}
\item $D_N(t)=\frac{\sin((N+\frac12)t)}{\sin(t/2)}$ for $t\notin2\pi\Z$, with $D_N(0)=2N+1$, and $s_N(f;x)=\frac{1}{2\pi}\int_{-\pi}^{\pi}f(x-t)D_N(t)\,dt$.
\item If $x\in\R$ and there exist $M,\delta>0$ such that $\abs{f(x+t)-f(x)}\leq M\abs{t}$ for $\abs{t}<\delta$, then $s_N(f;x)\to f(x)$.
\item If $g$ is Riemann integrable and $2\pi$-periodic and $f=g$ near $x$, then $s_N(f;x)-s_N(g;x)\to0$.
\end{enumerate}
\begin{proof}
\begin{enumerate}[label=(\roman*)]
\setlength{\itemsep}{0pt}
\item For $z=e^{it}\neq1$, $D_N(t)=z^{-N}\frac{1-z^{2N+1}}{1-z}=\frac{\sin((N+\frac12)t)}{\sin(t/2)}$, and $D_N(0)=\sum\limits_{m=-N}^{N}1=2N+1$.\\
Substituting $\widehat f(m)$ and interchanging the finite sum with the integral gives\\
$s_N(f;x)=\frac{1}{2\pi}\int_{-\pi}^{\pi}f(u)\sum\limits_{m=-N}^{N}e^{im(x-u)}\,du=\frac{1}{2\pi}\int_{-\pi}^{\pi}f(x-t)D_N(t)\,dt$.
\item Since $\frac{1}{2\pi}\int_{-\pi}^{\pi}D_N(t)\,dt=1$, (i) gives $s_N(f;x)-f(x)=\frac{1}{2\pi}\int_{-\pi}^{\pi}(f(x-t)-f(x))D_N(t)\,dt$.\\
Set $q_x(t)\defeq\frac{f(x-t)-f(x)}{\sin(t/2)}$ for $0<\abs{t}\leq\pi$ and $q_x(0)\defeq0$. The hypothesis makes $q_x$ bounded near $0$ and Riemann integrable on $[-\pi,\pi]$. Hence $s_N(f;x)-f(x)=\frac{1}{2\pi}\int_{-\pi}^{\pi}q_x(t)\sin((N+\frac12)t)\,dt$.\\
Since $\sin((N+\frac12)t)=\sin Nt\cos(t/2)+\cos Nt\sin(t/2)$, Theorem \ref{thm:orthonormal-fourier-best-approximation-bessel} applied to $q_x(t)\cos(t/2)$ and $q_x(t)\sin(t/2)$ sends both coefficients to $0$.
\item Apply (ii) to $h\defeq f-g$. Then $h=0$ near $x$, so $s_N(h;x)\to h(x)=0$.\\
By linearity, $s_N(h;x)=s_N(f;x)-s_N(g;x)$.
\end{enumerate}
\hlt{Structural Commentary}: Dirichlet kernels express partial sums as oscillatory local averages.\\
\hlt{Key Insight}: Local regularity makes the oscillatory error vanish.
\end{proof}
\end{theorem}

\begin{corollary} \label{cor:trigonometric-polynomial-weierstrass} \hlt{Uniform Approximation by Trigonometric Polynomials}\\
Let $f\in C(\R;\C)$ be $2\pi$-periodic. For every $\varepsilon>0$, there is $P(x)=\sum\limits_{n=-N}^{N}c_ne^{inx}$ such that $\norm{f-P}_\infty<\varepsilon$.
\begin{proof}
Let $\mathbb{T}\defeq\{z\in\C:\abs{z}=1\}$ and define $F:\mathbb{T}\to\C$ by $F(e^{ix})\defeq f(x)$.\\
Periodicity makes $F$ well-defined and continuous.\\
The algebra $\mathcal{T}\defeq\operatorname{span}\{z^n:n\in\Z\}\subset C(\mathbb{T})$ is self-adjoint, separates points, and vanishes nowhere.\\
By Theorem \ref{thm:complex-stone-weierstrass}, choose $Q(z)=\sum\limits_{n=-N}^{N}c_nz^n\in\mathcal{T}$ with $\norm{F-Q}_{C(\mathbb{T})}<\varepsilon$.\\
Set $P(x)\defeq Q(e^{ix})=\sum\limits_{n=-N}^{N}c_ne^{inx}$; then $\norm{f-P}_\infty=\norm{F-Q}_{C(\mathbb{T})}<\varepsilon$.\\
\hlt{Structural Commentary}: Trigonometric polynomials are Laurent polynomials pulled back from $\mathbb{T}$.\\
\hlt{Key Insight}: Stone--Weierstrass on $\mathbb{T}$ gives uniform periodic approximation.
\end{proof}
\end{corollary}

\begin{theorem} \label{thm:parseval-fourier-series} \hlt{Parseval's Theorem for Trigonometric Fourier Series}\\
Let $f,g$ be Riemann integrable and $2\pi$-periodic.\\
Set $\langle u,v\rangle_2\defeq\frac{1}{2\pi}\int_{-\pi}^{\pi}u(x)\overline{v(x)}\,dx$, $\norm{u}_2^2\defeq\langle u,u\rangle_2$, and $s_N(f)\defeq\sum\limits_{\abs{n}\leq N}\widehat f(n)e^{inx}$. Then:
\begin{enumerate}[label=(\roman*)]
\setlength{\itemsep}{0pt}
\item $\norm{f-s_N(f)}_2\to0$.
\item $\langle f,g\rangle_2=\sum\limits_{n\in\Z}\widehat f(n)\overline{\widehat g(n)}$.
\item $\norm{f}_2^2=\sum\limits_{n\in\Z}\abs{\widehat f(n)}^2$.
\end{enumerate}
\begin{proof}
\begin{enumerate}[label=(\roman*)]
\setlength{\itemsep}{0pt}
\item Put $d_N(u)\defeq\norm{u-s_N(u)}_2$. If $u$ is continuous, choose a trigonometric polynomial $P$ with $\norm{u-P}_2<\varepsilon$; for $N\geq\deg P$, Theorem \ref{thm:orthonormal-fourier-best-approximation-bessel} gives $d_N(u)\leq\norm{u-P}_2<\varepsilon$, so $d_N(u)\to0$.\\
For Riemann integrable $f$, choose continuous $h$ with $\norm{f-h}_2<\varepsilon$.\\
Since $\norm{s_N(h-f)}_2\leq\norm{h-f}_2$, $d_N(f)\leq2\norm{f-h}_2+d_N(h)$; hence $\limsup\limits_{N\to\infty}d_N(f)\leq2\varepsilon$. Let $\varepsilon\downarrow0$.
\item By orthogonality, $\langle e^{inx},g\rangle_2=\overline{\widehat g(n)}$. Hence $\langle s_N(f),g\rangle_2=\sum\limits_{\abs{n}\leq N}\widehat f(n)\langle e^{inx},g\rangle_2=\sum\limits_{\abs{n}\leq N}\widehat f(n)\overline{\widehat g(n)}$.\\
By (i) and Cauchy--Schwarz, $\langle s_N(f),g\rangle_2\to\langle f,g\rangle_2$.
\item Set $g=f$ in (ii). Then $\widehat g(n)=\widehat f(n)$ and $\langle f,f\rangle_2=\norm{f}_2^2$, so $\norm{f}_2^2=\sum\limits_{n\in\Z}\abs{\widehat f(n)}^2$.
\end{enumerate}
\hlt{Structural Commentary}: Density turns finite projection identities into $L^2$ identities.\\
\hlt{Key Insight}: Partial Fourier projections converge in $L^2$.
\end{proof}
\end{theorem}

\begin{method} \label{meth:fourier-structure-proof-patterns} \hlt{Fourier Proof Patterns}
\begin{enumerate}[label=(\roman*)]
\setlength{\itemsep}{0pt}
\item \hlt{Characters}: use orthogonality of $e_n$ to compute coefficients and finite-dimensional projections.
\item \hlt{Kernels}: rewrite partial sums or averages as convolutions with Dirichlet or Fejer kernels.
\item \hlt{Approximation}: use Fejer positivity or Stone--Weierstrass to approximate continuous periodic functions by trigonometric polynomials.
\item \hlt{$L^2$ control}: project onto finite orthonormal spans, then pass to limits by density.
\item \hlt{Localization}: express the error as an oscillatory integral and use local regularity near the point.
\end{enumerate}
\hlt{Key Insight}: Fourier theory is finite-dimensional orthogonal projection plus an approximation theorem.
\end{method}
